\documentclass[12pt, a4paper]{report}
\usepackage[utf8]{inputenc}
\usepackage{graphicx}
\usepackage{geometry}
\usepackage{hyperref}
\usepackage{titlesec}
\usepackage{parskip}

% Geometry settings
\geometry{
    a4paper,
    left=25mm,
    right=25mm,
    top=25mm,
    bottom=25mm
}

% Hyperref settings
\hypersetup{
    colorlinks=true,
    linkcolor=black,
    filecolor=magenta,      
    urlcolor=blue,
    pdftitle={CPU Round Robin Scheduling Algorithm Visualizer},
    pdfauthor={Rubin Bhattarai, Adarsha Pant}
}

% Title Format
\titleformat{\chapter}[display]
  {\normalfont\bfseries}{}{0pt}{\Huge}

\begin{document}

% -------------------------------------------------------------------
% TITLE PAGE
% -------------------------------------------------------------------
\begin{titlepage}
    \begin{center}
        \vspace*{1cm}
        
        \Huge
        \textbf{CPU Round Robin Scheduling Algorithm Visualizer}
        
        \vspace{0.5cm}
        \LARGE
        Computer Graphics Mini-Project Report
        
        \vspace{1.5cm}
        
        \textbf{Submitted by:} \\
        Rubin Bhattarai (Roll No: 16) \\
        Adarsha Pant (Roll No: 46)
        
        \vspace{2cm}

        \textbf{Department of Computer Engineering} \\
        [Your University/College Name] \\
        [City, Country]
        
        \vspace{2cm}
        
        \today
        
    \end{center}
\end{titlepage}

% -------------------------------------------------------------------
% ABSTRACT
% -------------------------------------------------------------------
\begin{abstract}
CPU scheduling is a fundamental concept in operating systems, often taught using static tables which fail to convey the dynamic nature of time slicing and process states. The \textbf{CPU Round Robin Scheduling Algorithm Visualizer} addresses this educational gap by providing an interactive graphical simulation. Utilizing a conveyor belt metaphor, it visually demonstrates the movement of process blocks between the Ready Queue and the CPU. This project leverages computer graphics techniques such as smooth motion simulation, state-driven rendering, and real-time data visualization to create an engaging learning tool. Key features include a live procedural Gantt chart, interactive parameters for time quantum and speed, and distinct visual states for processes. The system is implemented using JavaScript and the p5.js library.
\end{abstract}

\tableofcontents

% -------------------------------------------------------------------
% CHAPTER 1: INTRODUCTION
% -------------------------------------------------------------------
\chapter{Introduction}

\section{Background}
CPU scheduling is often taught using static tables, which makes it difficult to understand how time slicing actually behaves during execution. Students frequently struggle to visualize the concurrent nature of process management and the precise timing of context switches.

\section{Problem Statement}
Traditional methods of teaching Round Robin scheduling lack the visual feedback necessary to understand the "flow" of execution. Concepts like the Ready Queue, CPU burst times, and time quantums remain abstract numbers rather than tangible entities.

\section{Solution Overview}
The CPU Round Robin Visualizer addresses this gap by turning abstract operating system concepts into an interactive graphical simulation. The tool uses a conveyor belt style metaphor where color-coded process blocks move between the Ready Queue and the CPU. This movement helps users visually grasp how processes are scheduled and preempted when the time quantum expires. In addition, the system generates a live Gantt chart during execution, allowing users to see how scheduling history develops over time as the simulation runs.

% -------------------------------------------------------------------
% CHAPTER 2: PROJECT FEATURES
% -------------------------------------------------------------------
\chapter{Project Features}

The visualizer includes several key features designed to enhance the educational value and user experience:

\begin{itemize}
    \item \textbf{Smooth Motion Simulation}: Processes move using fluid animations instead of abrupt jumps, creating a more realistic sense of execution flow. This helps in tracking the lifecycle of a specific process visually.
    
    \item \textbf{Live Procedural Gantt Chart}: A dynamic timeline that grows and updates in real time, with color-coding to represent CPU activity. This provides an immediate historical view of the schedule.
    
    \item \textbf{Interactive Parameters}: Users can modify the time quantum and simulation speed to study their effect on scheduling behavior. This interactivity allows for "what-if" analysis.
    
    \item \textbf{State Visualization}: Clear visual representation of process states such as Ready, Running, and Terminated using distinct colors and remaining time indicators. The visual feedback is instant and state-dependent.
\end{itemize}

% -------------------------------------------------------------------
% CHAPTER 3: COMPUTER GRAPHICS CONCEPTS
% -------------------------------------------------------------------
\chapter{Computer Graphics Concepts Implemented}

This project is not just a logic simulator but a demonstration of computer graphics principles applied to educational software.

\section{Linear Interpolation (LERP)}
LERP is used to calculate smooth transitions as process blocks move between the queue, CPU, and completed states. Instead of snapping instantly from one position to another, the blocks smoothly glide, making it easier for the eye to follow the process flow.

\section{Coordinate Space Mapping}
The application converts abstract time units into pixel widths so the Gantt chart scales correctly within the display window. This involves mapping the logical domain (time) to the physical range (screen pixels).

\section{State Driven Rendering}
Visual attributes such as position, color, and size are directly controlled by the current state of each process. The rendering loop checks the state of every object each frame and draws it accordingly, ensuring the visual representation is always in sync with the logical model.

\section{Procedural UI Generation}
Two-dimensional shapes and text elements are generated dynamically based on live scheduling data. The UI is not a static image but a procedurally constructed scene that adapts to the number of processes and their states.

% -------------------------------------------------------------------
% CHAPTER 4: IMPLEMENTATION DETAILS
% -------------------------------------------------------------------
\chapter{Implementation Details}

\section{Technology Stack}
\begin{itemize}
    \item \textbf{Programming Language}: JavaScript
    \item \textbf{Graphics Library}: p5.js
\end{itemize}

\section{Overview}
The project is built using the p5.js library, which provides a canvas-based rendering context. The main simulation loop handles the logic of the Round Robin algorithm (decrementing burst times, checking time quantums, rotating the queue) and updates the visual properties of the process objects. The draw loop then renders these objects to the canvas.

% -------------------------------------------------------------------
% CHAPTER 5: CONCLUSION
% -------------------------------------------------------------------
\chapter{Conclusion}

The CPU Round Robin Scheduling Algorithm Visualizer successfully demonstrates how computer graphics can be effectively used to teach complex abstract concepts. By creating a visual metaphor for the CPU scheduling process, the project makes learning more intuitive and engaging. The use of smooth animations, real-time feedback, and interactive elements transforms a static problem into a dynamic experimentation playground.

\end{document}
